%! Author = kristofferpaulsson
%! Date = 2024-08-03

% Preamble
\documentclass[10pt,a4paper,notitlepage]{article}

\usepackage[margin=2.5cm]{geometry}
\usepackage{fontspec}
\usepackage[none]{hyphenat}  % --- NO HYPHENATION

\usepackage[hyphens]{xurl}  % --- URL HYPHENATION
\usepackage{hyperref}  % --- URL

\usepackage{draftwatermark}  % --- DONT-CITE

\urlstyle{same}  % --- URL
\pagenumbering{gobble}
\DraftwatermarkOptions{fontsize=2em,angle=0,vpos=1cm,color=black,text=Do Not Cite or Circulate}  % --- DONT-CITE

\usepackage{float}  % --- TABLE
\usepackage{makecell}  % --- TABLE
\usepackage{enumitem}  % --- LIST


\fontspec{Tinos}[
    Path = ../../font/t/,
    Extension = .ttf,
    UprightFont = *-Regular,
    Scale=1,
    Ligatures=TeX
]
\fontspec{Inter}[
    Path = ./../font/i/,
    Extension = .ttf ,
    UprightFont = *-Regular ,
    Scale=1,
    Ligatures=TeX
]

% Commands
\newcommand{\tsc}[1]{\textsc{#1}}
\newcommand*{\grc}[1]{\fontspec{Inter}#1\rmfamily}
\newcommand{\trc}[1]{\textit{\fontspec{Tinos}#1}}
\newcommand{\linkfoot}[3]{\footnote{\emph{#1}, s.v. ``{#2},'' accessed \printdate{#3}.}}

% Document
\begin{document}

    \subsection*{Ancient Greek: Reflexive Pronouns in First and Second Person}
    \begin{itemize}[label={}, leftmargin=0]
        \item Kristoffer E.D.\ Paulsson (2026), Independent Scholar.
    \end{itemize}
    Reflexive pronouns are used when the subject and object of a verb or preposition refer to the same entity,
    indicating that the action reflects back on the subject.
    \begin{quote}
    % Exact citation start
    The reflexive pronouns (referring back to the subject of the sentence) are formed by
    compounding the stems of the personal pronouns with the oblique cases of \grc{αὐτός}. In the plural both pronouns
    are declined separately, but the third person has also the compounded form. The nominative is excluded by
    the meaning. There is no dual.
    % Exact citation stop
    (Smyth, 1956:93, \S329)
    \end{quote}
    Rydberg-Cox (n.d.) says: ``The reflexive pronouns are compounded of the stems of the personal pronouns and
    \grc{αὐτός} [intensive pronouns]. But in the plural the two pronouns are declined separately in the first and
    second persons.'' It further states: ``The reflexive pronouns refer to the subject of the clause in which
    they stand. Sometimes in a dependent clause they refer to the subject of the leading verb,-- \textit{i.e.} they
    are \textit{indirect} reflexives.''

    \begin{table}[H]
        \begin{tabular}{c|cccc}
             & \multicolumn{2}{c}{\tsc{1st}} & \multicolumn{2}{c}{\tsc{2nd}} \\
             & \tsc{Masc} & \tsc{Fem} & \tsc{Masc} & \tsc{Fem} \\

            \hline
            \emph{\tsc{Singular}} & & & & \\
             \tsc{Nom} & \makecell{--} & \makecell{--}
            & \makecell{--} & \makecell{--}\\

            \tsc{Gen} & \makecell{\grc{ἐμαυτοῦ} \trc{emautou} \\ \small of myself (m.)} & \makecell{\grc{ἐμαυτῆς} \trc{emautēs} \\ \small of myself (f.)}
            & \makecell{\grc{σεαυτοῦ} \trc{seautou} \\ \small of yourself (m.)} & \makecell{\grc{σεαυτῆς} \trc{seautēs} \\ \small of yourself (f.)}\\

            \tsc{Dat} & \makecell{\grc{ἐμαυτῷ} \trc{emautōi} \\ \small to/for myself (m.)} & \makecell{\grc{ἐμαυτῇ} \trc{emautēi} \\ \small to/for myself (f.)}
            & \makecell{\grc{σεαυτῷ} \trc{seautōi} \\ \small to/for yourself (m.)} & \makecell{\grc{σεαυτῇ} \trc{seautēi} \\ \small to/for yourself (f.)}\\

            \tsc{Acc} & \makecell{\grc{ἐμαυτόν} \trc{emauton} \\ \small myself (m.)} & \makecell{\grc{ἐμαυτήν} \trc{emautēn} \\ \small myself (f.)}
            & \makecell{\grc{σεαυτόν} \trc{seauton} \\ \small yourself (m.)} & \makecell{\grc{σεαυτήν} \trc{seautēn} \\ \small yourself (f.)}\\

            \hline
            \emph{\tsc{Dual}} & \multicolumn{2}{c}{--} & \multicolumn{2}{c}{--} \\

            \hline
            \emph{\tsc{Plural}} & & & & \\
            \tsc{Nom} & \makecell{--} & \makecell{--}
            & \makecell{--} & \makecell{--}\\

            \tsc{Gen} & \multicolumn{2}{c}{\makecell{\grc{ἡμῶν αὐτῶν} \\ \trc{hēmōn autōn} \\ \small of ourselves}}
            & \multicolumn{2}{c}{\makecell{\grc{ὑμῶν αὐτῶν} \\ \trc{humōn autōn} \\ \small of yourselves}} \\

            %\emph{\tsc{}} & We/us men & We/us women & Yous men & Yous women \\
            \tsc{Dat} & \makecell{\grc{ἡμῖν αὐτοῖς} \\ \trc{hēmin autois} \\ \small to/for ourselves (m.)} & \makecell{\grc{ἡμῖν αὐταῖς} \\ \trc{hēmin autais} \\ \small to/for ourselves (f.)}
            & \makecell{\grc{ὑμῖν αὐτοῖς} \\ \trc{humin autois} \\ \small to/for yourselves (m.)} & \makecell{\grc{ὑμῖν αὐταῖς} \\ \trc{humin autais} \\ \small to/for yourselves (f.)}\\

            \tsc{Acc} & \makecell{\grc{ἡμᾶς αὐτούς} \\ \trc{hēmas autous} \\ \small ourselves (m.)} & \makecell{\grc{ἡμᾶς αὐτάς} \\ \trc{hēmas autas} \\ \small ourselves (f.)}
            & \makecell{\grc{ὑμᾶς αὐτούς} \\ \trc{humas autous} \\ \small yourselves (m.)} & \makecell{\grc{ὑμᾶς αὐτάς} \\ \trc{humas autas} \\ \small yourselves (f.)}\\
        \end{tabular}
        \caption{Compiled from standard Ancient Greek grammars enhanced with transliteration and cues. (Smyth, 1956:93, \S329)}
    \end{table}

    \subsection*{Reference List}
    \begin{itemize}[label={}, leftmargin=0]
        \item Rydberg-Cox, J. (n.d.). \textit{LESSON XLVIII: Reflexive, Reciprocal, and Possessive Pronouns}. [online] A Digital Tutorial for Ancient Greek. Available at: \url{https://daedalus.umkc.edu/FirstGreekBook/JWW_FGB48.html} [Accessed 26 Mar. 2026].
        %\item Cambridge Dictionary (2024). \textit{Yous}. [online] @CambridgeWords. Available at: \url{https://dictionary.cambridge.org/dictionary/english/yous} [Accessed 13 May 2025].
        \item Smyth, H.W.\ (1956). \textit{Greek Grammar}. Cambridge, Massachusetts: Harvard University Press.
    \end{itemize}

\end{document}